% this file is called up by thesis.tex
% content in this file will be fed into the main document

%: ----------------------- name of chapter  -------------------------

\chapter{Methodology}
\label{chap:methodology}

% change according to folder and file names
\ifpdf
    \graphicspath{{2_related_works/figures/PNG/}{2_related_works/figures/PDF/}{3_geprng/figures/}}
\else
    \graphicspath{{2_related_works/figures/EPS/}{2_related_works/figures/}}
\fi

%: ----------------------- contents from here ------------------------

This chapter presents the methodology used to explore the capacity of Deep Reconstructive Normal Behaviour Models in wind turbine condition monitoring. The proposed pipeline is designed to detect abnormal conditions based on telemetry data from real-world operational turbines. The methodology includes data preprocessing, model design, training setup, evaluation strategy, and deployment considerations.


% === Data Acquisition and Preprocessing ===
\section{Dataset}
\label{sec:dataset}

The proprietary dataset used in this research was provided by EnergyPro, an industrial partner aiming to bridge academic research and industrial applications. 

The dataset used in this study initially comprised 33 variables collected from operational wind turbines, encompassing a mix of physical sensor readings, control parameters, and static metadata. 
These features included both time-varying telemetry data, such as rotational speed and power output, as well as categorical and static attributes, including turbine site, model, and geographic coordinates.

To focus the modeling efforts on the most informative signals, a feature selection process was carried out in consultation with a domain expert. Based on the input, and further guided by theoretical and empirical considerations, six features were selected for inclusion in the final model: average power output (\texttt{avgpower}), average rotor speed (\texttt{avgrotorspeed}), average wind speed (\texttt{avgwindspeed}), air density (\texttt{density}), ambient temperature (\texttt{ambienttemperature}), and average wind direction (\texttt{avgwinddirection}).

This selection is not only consistent with expert judgment but is also grounded in fundamental aerodynamic theory. As described by \cite{bilendo_normal_2022}, the power generated by a wind turbine can be expressed by the following equation:


\[
P = \frac{1}{2} \rho A C_P(\lambda, \beta) v^3,
\]

where \(P\) denotes power output, \(\rho\) is air density, \(v\) is wind speed, and \(C_P(\lambda, \beta)\) is the power coefficient, which depends on the tip speed ratio \(\lambda\) and blade pitch angle \(\beta\). Within this formulation, at least four of the selected features—namely, wind speed, air density, rotor speed (which influences \(\lambda\)), and power output—are explicitly involved, while ambient temperature indirectly affects air density. This theoretical linkage reinforces their relevance for modeling turbine behaviour.


The remaining 27 features were excluded from further analysis for several reasons. First, a substantial portion of them consist of metadata—such as site identifier, manufacturer, model, latitude, and phase configuration—which do not change over time and are therefore uninformative for condition monitoring. Second, several features exhibited low signal variability or failed to show meaningful correlation with known abnormal events during exploratory analysis.

Focusing on a smaller, well-justified subset of features not only reduces the dimensionality of the input space but also mitigates the risk of overfitting and improves the interpretability of the learned models. This selection forms the basis for all subsequent preprocessing and modeling steps in this study.


% \section{Preprocessing}
% \begin{itemize}
%     \item \textit{Temporal Alignment:} Missing values interpolated using Pandas (\texttt{method='time'}), then averaged hourly.
%     \item \textit{Windowing:} Segmented into 128-hour sliding windows. Windows with $\geq$2 abnormal hours labeled as fully abnormal; others as normal.
%     \item \textit{Normalization:} Min-max scaled to EnergyPro's operational ranges: 
%     \begin{equation}
%         X_{\text{norm}} = \frac{X - X_{\min}}{X_{\max} - X_{\min}}
%     \end{equation}
%     \item \textit{Splits:}
%     \begin{itemize}
%         \item Unsupervised: Normal data randomly split 70\%/15\%/15\% (train/val/test).
%         \item Semi-supervised: Abnormal data split into \texttt{abn\_train}, \texttt{abn\_val}, and \texttt{abn\_test} for classifier head training.
%     \end{itemize}
% \end{itemize}

% % === Model Design ===
% \section{Model Design}
% \label{sec:model}

% Three autoencoder architectures were compared:
% \begin{itemize}
%     \item \textbf{CNN-AE:} Convolution layers for spatial feature extraction.
%     \item \textbf{VAE:} Probabilistic latent space with KL-divergence regularization.
%     \item \textbf{SAE:} Stacked dense layers with sparsity constraints.
% \end{itemize}

% \textbf{Variants Tested:}
% \begin{itemize}
%     \item Multihead attention vs. single-head.
%     \item Decoder ablation (full AE vs. encoder-only classifier).
% \end{itemize}

% \textbf{Loss Functions:}
% \begin{itemize}
%     \item \textit{MSE:} Standard reconstruction error:
%     \begin{equation}
%         \mathcal{L}_{\text{MSE}} = \frac{1}{n}\sum_{i=1}^n (x_i - \hat{x}_i)^2
%     \end{equation}
%     \item \textit{Fourier Loss:} Frequency-domain discrepancy (see Appendix~\ref{app:fourier}).
%     \item \textit{MMD:} Maximum Mean Discrepancy for latent-space normality.
% \end{itemize}

% % === Training Protocol ===
% \section{Training Protocol}
% \label{sec:training}

% \begin{itemize}
%     \item \textit{Unsupervised Phase:} Models trained on normal data only.
%     \item \textit{Semi-supervised Phase:} Classifier head fine-tuned with labeled abnormal data.
%     \item \textit{Hyperparameter:} Fixed across models:
%     \begin{itemize}
%         \item Latent dimension: 64
%         \item Batch size: 32
%         \item Optimizer: Adam ($\eta=10^{-3}$)
%     \end{itemize}
% \end{itemize}

% % === Evaluation Framework ===
% \section{Evaluation Framework}
% \label{sec:eval}

% \textbf{Metrics:}
% \begin{itemize}
%     \item Primary: AUC-ROC (threshold-agnostic).
%     \item Secondary: Accuracy, precision, recall, F1 (using Youden's threshold):
%     \begin{equation}
%         \text{Threshold}^* = \text{argmax}_{\tau} (\text{Recall}(\tau) + \text{Precision}(\tau) - 1)
%     \end{equation}
% \end{itemize}

% \textbf{Validation:}
% \begin{itemize}
%     \item Normal test data $\rightarrow$ reconstruction error distribution.
%     \item Abnormal test data $\rightarrow$ detection performance.
% \end{itemize}

% % === Reproducibility ===
% \section{Reproducibility}
% \label{sec:repro}

% \begin{itemize}
%     \item \textit{Code:} PyTorch \texttt{vX.Y.Z} (see Appendix~\ref{app:code}).
%     \item \textit{Hardware:} NVIDIA V100 GPU (16GB VRAM).
%     \item \textit{Randomness Control:} Seed fixed to \texttt{42}.
% \end{itemize}

